\conc{

This dissertation addressed the scientific problem of developing a methodology for teaching digital history to future professionals. The following main results were obtained:

\begin{enumerate}

\item A PRISMA-guided systematic review \cite{Page2021} of digital history in pedagogical research was conducted, identifying only 6 studies that met rigorous inclusion criteria from an initial pool of 847 records across six databases \cite{Korniienko2025cte}. This finding establishes the nascent state of digital history pedagogy as a research field and highlights the urgent need for empirical studies. The scarcity of research motivated the examination of gamification as the most extensively studied component of digital history methodology.

\item A comprehensive systematic review of gamification in history education was performed, synthesising evidence from 74 studies \cite{Korniienko2025els}. The review revealed: (a) severe definitional ambiguity, with only 14.9\% of studies providing an explicit definition of gamification \cite{Deterding2011}; (b) predominantly positive effects on engagement (70.3\% of studies) and learning outcomes (51.4\%), consistent with broader gamification meta-analyses \cite{Hamari2014, KoivistoHamari2019}; (c) six evidence-based design principles -- historical authenticity, narrative integration, progressive difficulty, immediate feedback, social elements, and multiple modalities; (d) four distinct conceptual approaches to gamification (structural, content, full, and game-based learning).

\item A methodological quality assessment of the 74 gamification studies was conducted using the MRQS framework \cite{Korniienko2025mrqs, Whiting2016}, revealing that 83.8\% exhibited a high risk of bias. The mean MRQS was 4.31/10 (SD = 1.93, range 0--9). The most critical weaknesses included small sample sizes (66.1\% with N$<$50), lack of control groups (only 31.1\% included comparison conditions), absence of operational definitions (85.1\%), and near-universal failure to report effect sizes (90.5\%). Risk-of-bias analysis identified participant selection (Domain~1, 85.1\% high risk) as the most persistent source of weakness. An LLM validation approach revealed that automated risk-of-bias assessment achieves only slight agreement with human reviewers (Cohen's $\kappa$ = $-$0.04 to 0.21) per the \citet{LandisKoch1977} interpretation scale, while factual data extraction is highly reliable (82--97\% agreement), establishing a hybrid model for educational research synthesis.

\item Implementation guidelines for gamification in history education were developed based on the synthesis of systematic reviews and alignment with the European Digital Competence Framework (DigComp~2.0) \cite{Vuorikari2016, Titova2025cte}. The alignment demonstrates that gamified history education simultaneously develops all five DigComp competence areas: information and data literacy, communication and collaboration, digital content creation, safety, and problem solving.

\item A six-month pedagogical experiment was designed and implemented at the Centre for Pre-University Training and Education of Kryvyi Rih State Pedagogical University (October 2024 -- April 2025). The course served 23 students and comprised 38 lectures, 30 seminars with 156 Google Forms quizzes, 7 modular controls, and the interactive gamified platform HistoryLikeGame with 23 game modules \cite{Korniienko2025aredu, Zamkova2023299}. This represents one of the most comprehensively documented gamification studies in history education to date \cite{Korniienko2025els}.

\item A Python-based data analysis pipeline was developed to systematically process all collected data, including chat log parsing, assessment timeline extraction, response analysis, attendance tracking, gamification cataloguing, and figure generation. The pipeline ensures reproducibility and transparency of the analytical process.

\item The findings were validated through cross-domain comparison with a parallel systematic review of gamification in software engineering education (29 studies) \cite{Korniienko2025aredu, Souza2018201}, which confirmed that the methodological challenges identified in history education are characteristic of gamification research across educational domains, including predominantly positive but uncertain effects (Cohen's d = 0.23 for learning outcomes), small sample sizes, and reliance on unvalidated instruments \cite{DichevDicheva2017}.

\item Analysis of the expanded response data (36 seminars spanning the full course, N=1--27 per seminar, plus 7 modular controls with N=7--27) revealed mean scores ranging from 33.3\% to 94.2\% across seminars and 63.8\% to 84.6\% across modular controls. The pre-gamification period (Seminars 1--24) showed a mean of 70.7\% across 23 seminars, while the post-gamification period (Seminars 25--39) showed 63.4\% across 13 seminars, though this difference is confounded by topic difficulty and assessment format. Over 200 individual data points across 90 named participants were tracked, enabling trajectory analysis. Four distinct performance profiles were identified: consistent high performers (mean 86.9\%), variable performers with emerging competence, struggling learners showing improvement through repeated attempts, and declining performers. The modular control data revealed a participation decline from 27 (MC~2) to 7--9 (MCs~5--7), with performance varying by content period.

\item The study provides unique documentation of education under wartime conditions. Attendance analysis identified 33 absence reports, with infrastructure-related causes (internet problems: 24\%; power outages: 9\%) directly attributable to the armed conflict. The multi-channel course design (Zoom, Viber, Google Drive, Google Forms, HistoryLikeGame) demonstrated resilience through redundancy, ensuring educational continuity despite infrastructure disruptions.

\end{enumerate}

\textbf{Scientific novelty.} The scientific novelty of the obtained results lies in the fact that:
\begin{itemize}
\item \textit{for the first time}, a comprehensive methodology for teaching digital history to future professionals integrating gamification \cite{Deterding2011} has been developed and theoretically substantiated; a PRISMA-guided systematic review of digital history pedagogy has been conducted \cite{Page2021, Korniienko2025cte}; the methodological quality of gamification research in history education has been assessed using MRQS with LLM validation \cite{Korniienko2025mrqs}; a gamification-enhanced history education experiment under wartime conditions has been documented \cite{Korniienko2025aredu};
\item \textit{improved}: the content and methods of history teaching using digital tools, through the development of an integrated assessment system and implementation guidelines based on DigComp~2.0 \cite{Vuorikari2016, Titova2025cte};
\item \textit{further developed}: the theoretical foundations for applying gamification and digital technologies in humanities education \cite{Korniienko2025els}, and the methodology for systematic reviews with automated quality assessment \cite{Korniienko2025mrqs}.
\end{itemize}

\textbf{Practical significance.} The practical results include: (1) the interactive gamified platform HistoryLikeGame; (2) a comprehensive set of 156 quiz forms for Ukrainian history NMT preparation; (3) the data analysis pipeline for processing educational data; (4) implementation guidelines for gamification in history education.

\textbf{Future research directions.}
\begin{enumerate}
\item \textit{Expanded data collection}: Collecting complete response data from all seminars and modular controls would enable before-after comparisons and more robust trajectory analysis.
\item \textit{Controlled study}: Implementing a quasi-experimental design with comparison groups (gamification vs. traditional) would strengthen causal inference \cite{Hattie2009, Sailer2017}.
\item \textit{Platform expansion}: Extending HistoryLikeGame to cover all historical periods and adding analytics tracking would provide more comprehensive gamification data.
\item \textit{Larger sample}: Replicating the study with multiple course sections would increase statistical power and generalisability \cite{Qian2016, KoivistoHamari2019, Sunkara2024203}.
\item \textit{LLM integration}: Exploring the use of LLMs not only for quality assessment \cite{Korniienko2025mrqs, Li2024} but also for adaptive quiz generation and personalised feedback.
\item \textit{Longitudinal tracking}: Following students through the NMT examination to assess the long-term impact of the preparatory course on examination performance.
\end{enumerate}

\textbf{Implementation.} The results of the dissertation have been implemented in the educational process of the Centre for Pre-University Training and Education of Kryvyi Rih State Pedagogical University (implementation certificate available upon request). The HistoryLikeGame platform continues to be publicly accessible for educational use.

.
}
